\section{问题重述与分析}
\label{sec:intro}

\subsection{问题背景}

微网是为解决光伏等分布式能源本地消纳、减少并网问题而发展起来的小型发用电系统。非孤岛式微网与外部电网（以下简称外网）相连：在外网电价较低或分布式能源有剩余电量的时段为储能设备充电，在外网电价较高的时段释放电能，从而在满足小区负载的前提下尽可能节省购电费用。因此，微网调控的核心是一类\textbf{带储能的时序能量调度问题}：利用小区负载、光伏发电量与储能设备当前储电量等数据，给出每个时段的计划购电量与储能充放电量。

附件1--附件4 给出了电价、小区负载、光伏发电预测与实际功率、光伏发电整点预报等数据；附录1 给出了储能设备的容量、功率与效率参数。四个问题在\textbf{信息结构}上逐步升级：问题1 为完全信息下的单日计划；问题2 中负载与光伏需要自行预测，且缺电须按 5 倍电价紧急购电；问题3 增加 6:00、12:00、18:00 三个时刻的光伏预报与付费调整机制；问题4 中电价本身也实时波动，需在未知电价的条件下重算问题2 与问题3。

\subsection{四问的信息结构与要求}

\cref{tab:problems} 对比了四问的信息结构与决策要求。

\begin{table}[H]
  \centering
  \caption{四个问题的信息结构与决策要求对比}
  \label{tab:problems}
  \small
  \begin{tabular}{p{1.4cm}p{4.0cm}p{3.4cm}p{5.4cm}}
    \toprule[1.5pt]
    问题 & 决策时可获得的信息 & 决策时点 & 目标与关键约束 \\
    \midrule[1pt]
    1 & 当日电价、小区负载、光伏预测功率全部已知（附件1） & 0:00 一次 & 全天购电费最小；供电不低于负载；储能 0:00 与 24:00 储电量相同 \\
    2 & 电价每天相同（附件1）；负载与光伏仅有历史数据（附件2），须自行预测 & 0:00 一次 & 计划购电费与紧急购电费之和的期望最小；缺电按交易时刻电价的 5 倍紧急购电 \\
    3 & 电价每天相同；负载须自预测；光伏有 0:00/6:00/12:00/18:00 的 24 小时整点预报（附件3） & 0:00、6:00、12:00、18:00 & 计划购电费、调整购电费、紧急购电费之和最小；调低部分按 0.5 倍、调高部分按 1.5 倍交易时刻电价计费 \\
    4 & 电价实时波动（附件4），0:00 决策时未知 & 同问题2、3 & 同问题2、3，但全部费用按实际交易时刻电价结算 \\
    \bottomrule[1.5pt]
  \end{tabular}
\end{table}

四个问题还要求按统一的表格格式输出：表1 给出指定时段的购电量与全天购电量、购电费；表2 给出储能设备在六个 4 小时时段的充放电量以及 0:00、24:00 的储电量；表3 给出指定日期（2025-03-20、06-21、09-23、12-21）的紧急购电量。问题3 另需分析\textbf{是否需要引入其他时刻的预报}制定调整购电策略。

\subsection{统一框架与核心量}

通读四问可以看到，尽管信息条件不同，其物理过程、几何约束与优化目标完全一致：\textbf{同一根母线上，购电量、光伏出力与储能放电之和等于负载与储能充电之和；储能电量按效率递推；以购电费用为目标}。据此本文建立全文统一框架：

\begin{quote}
\itshape 本文将所有问题统一归结为含储能的单母线能量流线性规划；四问的差别仅是"进入模型的序列是确定值、预测量，还是带更新的预测量"以及"偏差如何计价"。
\end{quote}

据此定义贯穿全文的三个核心量（\cref{tab:symbols}）：时段末储电量 $E_t$、计划购电量 $g_t$ 与净缺口 $n_t=(L_t-P_t)\Delta t$。四问递进关系为：确定信息下单日优化（问题1）$\to$ 预测替代真值并引入安全裕度（问题2）$\to$ 预报更新下的滚动重优化与偏差计价（问题3）$\to$ 电价亦不确定时按实际价格结算（问题4）。

\subsection{数据处理与关键统计}

\cref{tab:attachments} 汇总了四个附件的内容与规模。

\begin{table}[H]
  \centering
  \caption{附件数据内容与规模}
  \label{tab:attachments}
  \small
  \begin{tabular}{clc}
    \toprule[1.5pt]
    附件 & 内容 & 规模 \\
    \midrule[1pt]
    1 & 某天的电价、小区负载、光伏发电预测功率（10 分钟粒度） & $144\times4$ \\
    2 & 2025.1.1--12.31 每天的小区负载与光伏发电实际功率 & $2\times(365\times145)$ \\
    3 & 每天 0:00、6:00、12:00、18:00 发布的未来 24 小时整点光伏预报 & $1460\times26$ \\
    4 & 2025.1.1--12.31 每天的实时电价（10 分钟粒度） & $365\times145$ \\
    \bottomrule[1.5pt]
  \end{tabular}
\end{table}

数据处理上作了三项核对：其一，附件2、附件4 的时间标签为区间结束时刻（如 \texttt{00:10} 表示 $[0{:}00,0{:}10]$，\texttt{0:00+1} 表示 $[23{:}50,24{:}00]$），据此 144 个标签恰好完整覆盖 $[0{:}00,24{:}00]$，与"0:00 制定全天计划""0:00 与 24:00 储电量相同"的表述自洽；其二，附件3 的"预报 $k$ 小时"是\textbf{相对于发布时刻}的第 $k$ 小时（6:00 发布的"预报 1 小时"指当日 07:00），该对齐方式下各时段误差量级合理（见后文\cref{tab:fc-acc}），若误按"当日第 $k$ 小时"对齐将产生 6--18 小时的系统性错位；其三，附件2 存在格式差异（时间列为时间对象而非字符串），读取时统一转换为数值矩阵。

\begin{figure}[H]
  \centering
  \includegraphics[width=0.66\linewidth]{fig01.png}
  \caption{附件数据总览：(a) 全年小区负载热力图；(b) 全年光伏实际功率热力图；(c) 全年实时电价热力图；(d) 典型日（附件1）负载、光伏、净缺口与电价}
  \label{fig:data-overview}
\end{figure}

\cref{fig:data-overview} 与 \cref{tab:stats} 汇总了数据的主要特征：

\begin{table}[H]
  \centering
  \caption{关键统计量（实测）}
  \label{tab:stats}
  \small
  \begin{tabular}{lll}
    \toprule[1.5pt]
    对象 & 统计量 & 数值 \\
    \midrule[1pt]
    附件1 电价 & 最小值 / 最大值 / 均值 & $0.3713$ / $1.3952$ / $0.7662$ 元/kWh \\
    附件1 负载 & 最小值 / 最大值 / 均值 / 日电量 & $3309$ / $5959$ / $4626$ \si{kW} / $111025$ \si{kWh} \\
    附件1 光伏预测 & 峰值 / 日电量 / 占负载比例 & $7612$ \si{kW} / $55483$ \si{kWh} / $50.0\%$ \\
    附件2 负载 & 全年最小 / 最大 / 均值 & $1996$ / $7979$ / $4626$ \si{kW} \\
    附件2 光伏 & 日电量最小 / 均值 / 最大 & $34361$ / $55483$ / $75223$ \si{kWh} \\
    月度净缺口 & 12 月最大 / 5 月最小 & $76259$ / $38284$ kWh/日 \\
    附件4 电价 & 最小 / 最大 / 均值 & $0.0076$ / $1.7936$ / $0.7662$ 元/kWh \\
    附件4 日内极差 & 均值 / 最大 & $1.1073$ / $1.4995$ 元/kWh \\
    相关性 & 电价--负载 / 电价--光伏 & $0.508$ / $0.028$ \\
    \bottomrule[1.5pt]
  \end{tabular}
\end{table}

\cref{fig:monthly} 进一步给出月度特征。由 \cref{fig:data-overview,fig:monthly,tab:stats} 可得三点对建模直接有用的认识：

（1）\textbf{附件1 是全年典型日}。其负载与光伏曲线与附件2 的全年均值曲线相关系数均为 $1.0000$，说明附件1 并非任意一天，而是官方给定的"典型日"，可直接作为问题2 的基准预测轮廓或与滑动平均预测作对照。

（2）\textbf{储能容量相对缺口偏小}。可用容量为 $10800-1200=9600$ \si{kWh}，仅占典型日净缺口（$55542$ \si{kWh}）的 $17\%$、日均负载电量的 $8.6\%$。因此储能的角色是\textbf{时段间的价格搬运}而非替代购电，全天购电量不可能低于净缺口。

（3）\textbf{季节差异显著}。12 月与 1 月的日均净缺口（$76259$、$75278$ \si{kWh}）约为 5 月（$38284$ \si{kWh}）的 2 倍，且 12 月电价最高（月均 $0.8617$ 元/kWh）；光伏日电量在 12 月仅 $36332$ \si{kWh}、8 月可达 $69984$ \si{kWh}。这意味着预测误差与紧急购电风险在冬季最集中，安全裕度应随季节调整。

\begin{figure}[H]
  \centering
  \includegraphics[width=0.66\linewidth]{fig02.png}
  \caption{月度特征：(a) 光伏日电量与净缺口电量的月度分布；(b) 日均负载与月均电价的月度特征}
  \label{fig:monthly}
\end{figure}

\begin{figure}[H]
  \centering
  \includegraphics[width=0.66\linewidth]{fig03.png}
  \caption{表3 指定四天（2025-03-20、06-21、09-23、12-21）的负载、光伏与实时电价}
  \label{fig:fourdays}
\end{figure}

\cref{fig:fourdays} 给出了表3 四个指定日期的数据画像：06-21（夏至）光伏日电量最高（$62073$ \si{kWh}）、净缺口最小（$19335$ \si{kWh}）；12-21（冬至）光伏仅 $35339$ \si{kWh}、净缺口达 $89038$ \si{kWh} 且电价均值最高（$0.9101$ 元/kWh），是最严酷的调度场景。这四天恰好覆盖光伏出力的四季极端，后文各问均按此四天给出明细结果。
